% Luc Destombe --- licence CC BY-NC-SA 4.0
% https://creativecommons.org/licenses/by-nc-sa/4.0/deed.fr
% Fichier autonome : compiler avec pdflatex, deux fois (signets, nombre de pages).
\documentclass[a4paper,10pt]{article}
\makeatletter
\newif\iflycee@niveaux
\newif\iflycee@palatino
\lycee@palatinotrue

\RequirePackage[utf8]{inputenc}
\RequirePackage[T1]{fontenc}
\RequirePackage[french]{babel}
\RequirePackage{etoolbox}

\newcommand{\ShorthandsOff}{\@ifundefined{shorthandoff}{}{\shorthandoff{;:!?}}}
\newcommand{\ShorthandsOn}{\@ifundefined{shorthandon}{}{\shorthandon{;:!?}}}
\ShorthandsOff
\AtBeginDocument{\ShorthandsOn}
\AtBeginEnvironment{tikzpicture}{\ShorthandsOff}

\RequirePackage{geometry}
\RequirePackage{fancyhdr}
\RequirePackage{lastpage}
\RequirePackage{multicol}
\RequirePackage{array}
\RequirePackage{enumitem}
\RequirePackage{xcolor}
\RequirePackage{needspace}

\RequirePackage{amsmath,amssymb}
\RequirePackage{mathpazo}
\DeclareMathSizes{10.95}{10.95}{8}{5.5}
\begingroup\catcode`\_=\active \catcode`\^=\active
\gdef_#1{\sb{\scriptscriptstyle #1}}
\gdef^#1{\sp{\scriptscriptstyle\lycee@moinsepais #1}}
\endgroup
\newcommand\lycee@moins{\mathbin{%
  \setbox\z@\hbox{$\m@th\scriptscriptstyle\lycee@moinsorig$}%
  \dimen@=.55\wd\z@ \advance\dimen@-.5pt
  \hbox to.75\wd\z@{\hss$\m@th\scriptscriptstyle\vcenter{\hbox{%
    \tikz\draw[line width=.5pt,line cap=round](0,0)--(\the\dimen@,0);}}$\hss}}}
\begingroup\lccode`\~=`\- \lowercase{\endgroup\let~}\lycee@moins
\newcommand\lycee@moinsepais{\mathcode`\-="8000 }
\AtBeginDocument{\mathchardef\lycee@moinsorig=\mathcode`\-\relax}
\AtBeginDocument{\catcode`\_=12 \mathcode`\_="8000
  \catcode`\^=12 \mathcode`\^="8000 }
\newcommand\lycee@exposants{%
  \fontdimen13\textfont2=.48\fontdimen6\textfont2
  \fontdimen14\textfont2=.48\fontdimen6\textfont2
  \fontdimen15\textfont2=.48\fontdimen6\textfont2
  \fontdimen13\scriptfont2=.48\fontdimen6\scriptfont2
  \fontdimen14\scriptfont2=.48\fontdimen6\scriptfont2
  \fontdimen15\scriptfont2=.48\fontdimen6\scriptfont2}
\every@math@size\expandafter{\the\every@math@size\lycee@exposants}
\newlength{\retraitcalculs}
\setlength{\retraitcalculs}{2em}

\RequirePackage{tikz}
\usetikzlibrary{arrows.meta,calc,babel,shapes.misc}
\RequirePackage[most]{tcolorbox}
\tcbuselibrary{skins,breakable}

\definecolor{coulDef}{HTML}{1F4E79}
\definecolor{coulProp}{HTML}{7030A0}
\definecolor{coulRem}{HTML}{C55A11}
\definecolor{coulEx}{HTML}{2E7D32}
\definecolor{coulExo}{HTML}{B03A2E}
\definecolor{coulPart}{HTML}{1F4E79}
\definecolor{coulMeth}{HTML}{1B6E4F}
\definecolor{coulCourbe}{HTML}{0B5ED7}
\definecolor{coulCourbeB}{HTML}{C00000}
\definecolor{coulCourbeC}{HTML}{2E7D32}
\definecolor{coulGrille}{HTML}{8C8C8C}
\definecolor{coulRep}{HTML}{1B6E4F}
\definecolor{coulBar}{HTML}{2E7D32}

\newcommand{\R}{\mathbb{R}}
\newcommand{\Rs}{\mathbb{R}^{*}}
\renewcommand{\le}{\leqslant}
\renewcommand{\ge}{\geqslant}
\newcommand{\vect}[1]{\overrightarrow{#1}}
\newcommand{\norme}[1]{\left\lVert #1 \right\rVert}
\newcommand{\intervalle}[2]{\left[#1\,;#2\right]}
\RequirePackage{cancel}
\renewcommand{\CancelColor}{\color{coulExo}}
\newcommand{\fois}[1]{\times{\color{coulExo}#1}}
\newcommand{\barre}[1]{\cancel{#1}}

\tikzset{grille/.style={coulGrille,line width=0.4pt}}
\newcommand{\quadrillage}[4]{\draw[grille,step=1] (#1,#2) grid (#3,#4);}
\tikzset{
  croix/.style={cross out,draw,line width=0.6pt,inner sep=0pt,minimum size=4pt},
  croixdroite/.style={croix,rotate=45,line width=0.8pt},
}
\newcommand{\pt}[3]{\path (#1) node[croix]{} node[#2,font=\small,inner sep=2.5pt]{$#3$};}

\newcommand{\lycee@repere}[4]{%
  \draw[grille,step=1] (#1,#2) grid (#3,#4);
  \draw[-{Stealth[length=2mm]},thick] (#1,0)--({#3+0.4},0) node[below left=-1pt]{$x$};
  \draw[-{Stealth[length=2mm]},thick] (0,#2)--(0,{#4+0.4}) node[below left=-1pt]{$y$};
  \node[below left,font=\scriptsize,inner sep=1.5pt] at (0,0) {$0$};}
\newcommand{\lycee@gradx}[1]{\foreach \k/\l in {#1}{%
  \draw (\k,0.07)--(\k,-0.07) node[below,font=\scriptsize,inner sep=1.5pt]{$\l$};}}
\newcommand{\lycee@grady}[1]{\foreach \k/\l in {#1}{%
  \draw (0.07,\k)--(-0.07,\k) node[left,font=\scriptsize,inner sep=1.5pt]{$\l$};}}
\newcommand{\lycee@intff}[2]{\left[#1\,;#2\right]}
\newcommand{\lycee@intfo}[2]{\left[#1\,;#2\right[}
\newcommand{\lycee@intof}[2]{\left]#1\,;#2\right]}
\newcommand{\lycee@intoo}[2]{\left]#1\,;#2\right[}
\newcommand{\lycee@ens}[1]{\left\{#1\right\}}
\AtBeginDocument{%
  \providecommand{\repere}{\lycee@repere}%
  \providecommand{\gradx}{\lycee@gradx}%
  \providecommand{\grady}{\lycee@grady}%
  \providecommand{\Cf}{\mathcal{C}\sb{\scriptscriptstyle f}}%
  \providecommand{\Cg}{\mathcal{C}\sb{\scriptscriptstyle g}}%
  \providecommand{\intff}{\lycee@intff}%
  \providecommand{\intfo}{\lycee@intfo}%
  \providecommand{\intof}{\lycee@intof}%
  \providecommand{\intoo}{\lycee@intoo}%
  \providecommand{\ens}{\lycee@ens}}

\newlist{questions}{enumerate}{3}
\setlist[questions,1]{label=\arabic*\textdegree),leftmargin=*,itemsep=2pt,topsep=3pt}
\setlist[questions,2]{label=\alph*),leftmargin=*,itemsep=1pt,topsep=2pt}
\setlist[questions,3]{label=\roman*),leftmargin=*,itemsep=1pt,topsep=1pt}
\setlist[itemize]{label=\textbullet,leftmargin=*,itemsep=2pt,topsep=3pt}

\newcommand{\besoinplace}[1]{\par\penalty\@M\begingroup
  \dimen@=#1\relax\dimen@ii=\pagegoal\advance\dimen@ii-\pagetotal
  \ifdim\dimen@>\dimen@ii\ifdim\dimen@ii>\z@\vfil\fi\break\fi\endgroup}

\newcommand{\lycee@pied}{\thepage/\pageref{LastPage}}
\newcommand{\entete}[2]{%
  \pagestyle{fancy}\fancyhf{}%
  \fancyhead[L]{\footnotesize\color{coulPart}\textbf{#1}}%
  \fancyhead[R]{\footnotesize\color{coulPart}\textbf{#2}}%
  \fancyfoot[C]{\footnotesize\lycee@pied}%
  \renewcommand{\headrulewidth}{0.4pt}%
  \renewcommand{\headrule}{\hbox to\headwidth{%
    \color{coulPart!40}\leaders\hrule height \headrulewidth\hfill}}}

\RequirePackage{ccicons}
\newcommand{\lycee@licence}{{\scriptsize\color{gray}%
  \copyright\ Luc Destombe --- \ccbyncsa\ CC BY-NC-SA 4.0}}
\newif\iflycee@licence \lycee@licencetrue
\newcommand{\sanslicence}{\lycee@licencefalse}
\AtBeginDocument{\iflycee@licence\AddToHookNext{shipout/foreground}{%
  \put(0,0){\rlap{\kern\dimexpr1in+\hoffset+\oddsidemargin+\textwidth\relax
    \llap{\raisebox{-\dimexpr1in+\voffset+\topmargin+\headheight+\headsep
      +\textheight+\footskip\relax}{\lycee@licence}}}}}\fi}

\newcommand{\formulesserrees}{%
  \setlength{\abovedisplayskip}{3pt}\setlength{\belowdisplayskip}{3pt}%
  \setlength{\abovedisplayshortskip}{2pt}\setlength{\belowdisplayshortskip}{3pt}}

\setlength{\parindent}{0pt}

\RequirePackage[implicit=false,hidelinks,pdfencoding=auto,psdextra,bookmarksopen,
  bookmarksopenlevel=1,pdfcreator={},pdfproducer={}]{hyperref}
\RequirePackage{bookmark}
\hypersetup{pdfauthor={Luc Destombe},pdfkeywords={CC BY-NC-SA 4.0}}
\newcounter{lycee@signet}
\newcommand{\signet}[3][]{\stepcounter{lycee@signet}%
  \edef\lycee@dest{\if\relax\detokenize{#1}\relax
    signet.\arabic{lycee@signet}\else#1\fi}%
  \hypertarget{\lycee@dest}{}\bookmark[level=#2,dest=\lycee@dest]{#3}}
\pdfstringdefDisableCommands{%
  \def\R{R}\def\vect#1{#1}\def\Cf{Cf}\def\Cg{Cg}\def\fois#1{×#1}%
  \def\barre#1{#1}\def\intervalle#1#2{[#1;#2]}\def\intff#1#2{[#1;#2]}%
  \def\textsuperscript#1{#1}\def\dfrac#1#2{#1/#2}\def\frac#1#2{#1/#2}%
  \def\vec#1{vecteur #1}}
\RequirePackage[official]{eurosym}

\geometry{top=0.8cm,bottom=1.3cm,left=1cm,right=1cm,footskip=0.6cm}

\pagestyle{fancy}
\fancyhf{}
\renewcommand{\headrulewidth}{0pt}
\fancyfoot[C]{\footnotesize Page \thepage{} / \pageref{LastPage}}

\newcommand{\titrefiche}[2]{%
  \begin{center}
    \Large\textbf{#1}\\[2pt]
    \large\textbf{#2}\\[2pt]
  \end{center}
  \vspace{0.10cm}}

\setlength{\columnseprule}{0.4pt}
\setlength{\columnsep}{15pt}
\renewcommand{\columnseprulecolor}{\color{black!45}}
\setlength{\multicolsep}{2pt}

\newcounter{partiefiche}
\renewcommand{\thepartiefiche}{\Roman{partiefiche}}
\newif\ifapresbandeau
\newcommand{\lycee@bandeau}[1]{%
  \par\penalty-600
  \vspace{0.08cm}%
  \noindent\colorbox{coulPart}{%
    \parbox{\dimexpr\linewidth-2\fboxsep\relax}{%
      \color{white}\bfseries\raggedright\strut#1\strut}}%
  \nopagebreak\par\nopagebreak
  \vspace{0.06cm}\nopagebreak
  \apresbandeautrue}
\newcommand{\partiefiche}[1]{%
  \refstepcounter{partiefiche}%
  \lycee@bandeau{\signet{1}{\thepartiefiche{} --- #1}\thepartiefiche{} --- #1}}
\newcommand{\partiefichesansnum}[1]{\lycee@bandeau{\signet{1}{#1}#1}}

\newcounter{exo}
\newlength{\espaceexo}\setlength{\espaceexo}{7pt}
\newlength{\espacetitre}\setlength{\espacetitre}{2pt}
\newcommand{\exo}[1][\relax]{%
  \par
  \ifapresbandeau\apresbandeaufalse\else\penalty-150\addvspace{\espaceexo}\fi
  \refstepcounter{exo}%
  \noindent\signet[exercice-\arabic{exo}]{2}{Exercice \arabic{exo}}%
  \textbf{\color{coulExo}\underline{Exercice \theexo}}%
  \ifx\relax#1\relax\else\textbf{ : #1}\fi
  \par\nopagebreak\vspace{\espacetitre}\nopagebreak
  \ignorespaces}

\setlist[enumerate]{nosep,leftmargin=*,itemsep=2pt,topsep=2pt}
\setlist[itemize]{nosep,leftmargin=*,topsep=2pt}
\newcommand{\choix}[3]{\par\hspace*{1em}%
  \makebox[0.3\linewidth][l]{a) #1}\makebox[0.3\linewidth][l]{b) #2}c) #3}
\newcommand{\deux}[2]{\par\hspace*{0.6em}%
  \makebox[0.48\linewidth][l]{#1}#2}
\newcommand{\trois}[3]{\par\hspace*{0.6em}%
  \makebox[0.32\linewidth][l]{#1}\makebox[0.32\linewidth][l]{#2}#3}

\renewenvironment{center}{\par\addvspace{3pt}\centering}{\par\addvspace{3pt}}
\formulesserrees
\AtBeginDocument{\formulesserrees}
\clubpenalty=10000
\widowpenalty=10000
\displaywidowpenalty=10000
\brokenpenalty=10000
\setlength{\parskip}{0pt}

\RequirePackage{ragged2e}
\setlength{\RaggedRightRightskip}{0pt plus 1fil}
\AtBeginDocument{\RaggedRight\binoppenalty=10000 \relpenalty=10000 }
\g@addto@macro\@parboxrestore{\RaggedRight}
\makeatother

\definecolor{coulInt}{HTML}{D62828}
\newcommand{\axegrad}[2]{%
  \draw[-{Stealth[length=2mm]}] ({#1-0.5},0)--({#2+0.5},0);
  \foreach \k in {#1,...,#2} \draw (\k,0.07)--(\k,-0.07);}

\begin{document}

\titrefiche{Seconde --- Intervalles et fonctions}{Fiche d'exercices du chapitre}

\begin{multicols}{2}\raggedcolumns

\partiefiche{Droite numérique, intervalles fermés}

\exo[Représenter des intervalles fermés]
\begin{enumerate}[label=\arabic*)]
  \item Représenter sur une droite graduée :\\
  $I=\intff{1}{4}$ ; \ $J=\intff{-6}{-3}$ ; \ $K=\intff{-2}{3}$ ; \
  $L=\intff{-\frac12}{\frac52}$.
  \item Écrire l'intervalle représenté ci-dessous.
  \begin{center}
  \begin{tikzpicture}[x=0.55cm]
    \axegrad{-4}{6}
    \node[below=6pt,font=\scriptsize] at (0,0) {$0$};
    \node[below=6pt,font=\scriptsize] at (1,0) {$1$};
    \draw[coulInt,line width=1.8pt] (-2,0)--(5,0);
    \draw[coulInt,line width=1.2pt] ($(-2,0)+(2.6pt,5pt)$) -- ++(-2.6pt,0) -- ++(0,-10pt) -- ++(2.6pt,0);
    \draw[coulInt,line width=1.2pt] ($(5,0)+(-2.6pt,5pt)$) -- ++(2.6pt,0) -- ++(0,-10pt) -- ++(-2.6pt,0);
  \end{tikzpicture}
  \end{center}
  \item Pour chaque nombre, dire s'il appartient à l'intervalle écrit à côté :\\
  a)~$4$ et $\intff{0}{4}$ ; \ b)~$9$ et $\intff{3}{8{,}9}$ ; \ c)~$-3$ et
  $\intff{-4}{-2}$ ; \ d)~$-50$ et $\intff{-49}{-5}$.
\end{enumerate}

\exo[Appartenance]
Compléter par $\in$ ou $\notin$ :
\begin{multicols}{2}
\begin{enumerate}[label=\alph*)]
  \item $2\;\ldots\;\intff{-3}{7}$
  \item $-4\;\ldots\;\intff{-3}{5}$
  \item $10^{-2}\;\ldots\;\intff{0}{1}$
  \item $\pi\;\ldots\;\intff{3{,}1}{3{,}2}$
  \item $-3\;\ldots\;\intff{-6}{-3}$
  \item $1{,}5\;\ldots\;\intff{1{,}6}{2}$
\end{enumerate}
\end{multicols}

\exo[Nombres dans un intervalle]
\begin{enumerate}[label=\arabic*)]
  \item Donner tous les entiers positifs ou nuls de l'intervalle $\intff{-2}{\frac72}$.
  \item Donner deux nombres décimaux de $\intff{3{,}7}{3{,}8}$, autres que les bornes.
  \item Donner un intervalle fermé qui contient $\pi$ et dont les bornes sont deux
  entiers qui se suivent.
\end{enumerate}

\partiefiche{Intervalles ouverts, inégalités}

\exo[D'un intervalle à une inégalité]
Écrire avec une inégalité le fait que $x$ appartient à chaque intervalle. Dire si
l'intervalle est ouvert, fermé ou semi-ouvert.
\begin{multicols}{2}
\begin{enumerate}[label=\alph*)]
  \item $\intff{2}{7}$
  \item $\intoo{-1{,}5}{+\infty}$
  \item $\intof{-4}{0}$
  \item $\intoo{-\infty}{5}$
  \item $\intfo{-3}{8}$
  \item $\intoo{0}{1}$
\end{enumerate}
\end{multicols}

\exo[D'une phrase à un intervalle]
Chaque phrase décrit des nombres $x$. Écrire l'inégalité qui les décrit, puis
l'intervalle, et le représenter sur une droite graduée.
\begin{enumerate}[label=\arabic*)]
  \item Les nombres supérieurs ou égaux à $-1$.
  \item Les nombres strictement compris entre $-3$ et $2$.
  \item Les nombres strictement inférieurs à $4$.
  \item Les nombres positifs ou nuls et strictement inférieurs à $6$.
\end{enumerate}

\exo[D'une inégalité à un intervalle]
Écrire avec un intervalle, puis représenter sur une droite graduée :
\begin{multicols}{3}
\begin{enumerate}[label=\alph*)]
  \item $-4\le x\le 2$
  \item $x<3$
  \item $x>5$
  \item $2<x<3$
  \item $x\ge 1{,}5$
  \item $x\le 10^{-2}$
\end{enumerate}
\end{multicols}

\exo[Tableau à compléter]
\begin{enumerate}[label=\arabic*)]
  \item Écrire avec un intervalle :
  \choix{$\frac14\le x<\frac23$}{$x>\pi$}{$x<-1$}
  \item Compléter le tableau (représentation sur une droite graduée à main levée).
\end{enumerate}
\begin{center}
\renewcommand{\arraystretch}{1.35}
\small
\begin{tabular}{|c|c|p{3.4cm}|}
\hline
\textbf{Inégalité(s)} & \textbf{Intervalle} & \textbf{Représentation}\\\hline
$-1\le x\le 4$ & & \\\hline
 & $\intof{1}{5}$ & \\\hline
$-3\le x<2$ & & \\\hline
$0<x<3$ & & \\\hline
$x>2$ & & \\\hline
$x\le 1$ & & \\\hline
 & $\intfo{-2}{6}$ & \\\hline
\end{tabular}
\end{center}

\exo[Bornes incluses ou exclues ?]
Compléter par $\in$ ou $\notin$, en s'aidant des crochets.
\begin{multicols}{2}
\begin{enumerate}[label=\alph*)]
  \item $6\;\ldots\;\intof{-2}{6}$
  \item $-2\;\ldots\;\intoo{-2}{4}$
  \item $10^{-2}\;\ldots\;\intoo{0}{+\infty}$
  \item $\pi\;\ldots\;\intoo{3{,}1}{3{,}2}$
  \item $-4\;\ldots\;\intoo{-\infty}{-4}$
  \item $-1{,}5\;\ldots\;\intfo{-1{,}5}{1{,}5}$
  \item $0\;\ldots\;\intoo{-\infty}{+\infty}$
  \item $-1\;\ldots\;\intof{-1}{3}$
\end{enumerate}
\end{multicols}

\partiefiche{Intersection, réunion, complémentaire}

\exo[Premières intersections et réunions]
Pour chaque couple d'intervalles $I$ et $J$, déterminer $I\cap J$ et $I\cup J$ (on peut
représenter $I$ et $J$ sur une même droite graduée).
\begin{enumerate}[label=\alph*)]
  \item $I=\intfo{-1{,}5}{3}$ et $J=\intof{0}{6}$
  \item $I=\intoo{-\infty}{3}$ et $J=\intfo{1}{+\infty}$
\end{enumerate}

\exo[Attention aux bornes]
Même consigne.
\begin{enumerate}[label=\alph*)]
  \item $I=\intfo{-2}{+\infty}$ et $J=\intof{-3}{1}$
  \item $I=\intfo{3}{6{,}5}$ et $J=\intoo{1}{3}$
  \item $I=\intof{-\infty}{2}$ et $J=\intoo{-\infty}{4}$
  \item $I=\intof{-\infty}{5}$ et $J=\intfo{5}{+\infty}$
  \item $I=\intfo{-1}{+\infty}$ et $J=\intoo{-\infty}{-1}$
\end{enumerate}

\exo[Plusieurs intervalles]
\begin{enumerate}[label=\arabic*)]
  \item Écrire avec un intervalle l'ensemble des nombres $x$ tels que :\\[1pt]
  \begin{tabular}{@{}l@{\qquad\qquad}l@{}}
    $I$ : $-2\le x\le 6$ & $L$ : $x\le 6$\\
    $J$ : $-4\le x<-1$ & $M$ : $2<x$
  \end{tabular}
  \item On pose aussi $K=\intof{1}{10}$. Déterminer :\\
  $I\cap J$ ; \quad $I\cup J$ ; \quad $J\cap K$ ; \quad $K\cup L$ ; \quad $I\cap L$ ; \quad $K\cup M$.
\end{enumerate}

\exo[Complémentaire]
Pour chaque ensemble, déterminer son complémentaire dans $\R$.
\begin{multicols}{2}
\begin{enumerate}[label=\alph*)]
  \item $A=\intoo{3}{+\infty}$
  \item $B=\intof{-\infty}{4}$
  \item $C=\intfo{-2}{+\infty}$
  \item $D=\intoo{-\infty}{6}$
  \item $E=\intff{1}{2}$
  \item $F=\intoo{-5}{8}$
\end{enumerate}
\end{multicols}

\exo[Inégalités et intervalles : « et », « ou »]
Compléter le tableau : sur chaque ligne, un même ensemble de nombres $x$ est décrit par
des inégalités, puis par des intervalles. Deux conditions reliées par \og et\fg{} donnent
une intersection ($\cap$), reliées par \og ou\fg{} une réunion ($\cup$). Écrire ensuite
le résultat le plus simplement possible quand on le peut.
\begin{center}
\renewcommand{\arraystretch}{1.45}
\small
\setlength{\tabcolsep}{4pt}
\begin{tabular}{|c|>{\centering\arraybackslash}p{4cm}|>{\centering\arraybackslash}p{3.2cm}|}
\hline
 & \textbf{Inégalités} & \textbf{Intervalles}\\\hline
a) & $-3<x\le -1$ \ ou \ $2\le x<5$ & \\\hline
b) & & $\intff{-2}{0}\cup\intfo{3}{+\infty}$\\\hline
c) & $x<-1$ \ ou \ $x>4$ & \\\hline
d) & & $\intof{-\infty}{-3}\cup\intoo{2}{+\infty}$\\\hline
e) & $x>-2$ \ et \ $x\le 3$ & \\\hline
f) & & $\intfo{1}{+\infty}\cap\intoo{-\infty}{5}$\\\hline
g) & $x\ge -4$ \ et \ $x\ge 1$ & \\\hline
h) & $x<0$ \ ou \ $x\le 2$ & \\\hline
i) & $x<-1$ \ et \ $x>3$ & \\\hline
\end{tabular}
\end{center}

\partiefiche{Image et antécédent}

\exo[Calculer des images]
On considère la fonction ${f : x\longmapsto 2x-3x^{2}}$.
\begin{enumerate}[label=\arabic*)]
  \item Calculer l'image par $f$ de :\choix{$1$}{$-2$}{$\frac12$}
  \item Le nombre $2$ est-il un antécédent de $-8$ par $f$ ? Justifier.
\end{enumerate}

\exo[Image et antécédent]
On considère la fonction $g$ définie par ${g(x)=4x-5}$.
\begin{enumerate}[label=\arabic*)]
  \item Calculer l'image par $g$ de :\choix{$-3$}{$5$}{$\frac34$}
  \item Le nombre $\frac12$ est-il un antécédent de $3$ par $g$ ? Justifier.
  \item Déterminer l'antécédent de $3$ par $g$.
\end{enumerate}

\exo[Calculer des antécédents]
Pour chaque fonction $f$, calculer le ou les antécédents du nombre $y$.
\begin{multicols}{2}
\begin{enumerate}[label=\alph*)]
  \item ${f(x)=x-4}$, \ $y=-2$
  \item ${f(x)=x^{2}}$, \ $y=36$
  \item ${f(x)=\frac1x}$, \ $y=4$
  \item ${f(x)=x(x-3)}$, \ $y=0$
\end{enumerate}
\end{multicols}

\exo[Deux fonctions égales]
On considère les fonctions $f$ et $g$ définies sur $\R$ par
\[
  f(x)=(x+4)(3-x)+x^{2}-10 \qquad\text{et}\qquad g(x)=2-x.
\]
\begin{enumerate}[label=\arabic*)]
  \item Calculer l'image de $2$ par $f$, puis par $g$.
  \item Choisir un nombre et calculer son image par $f$, puis par $g$.
  \item Démontrer que tout nombre a la même image par $f$ et par $g$.
\end{enumerate}

\partiefiche{Ensemble de définition}

\exo[Un nombre sans image]
Pour chaque fonction ci-dessous, un des nombres $-5$ ; $2$ ; $0$ ; $9$ n'a pas d'image.
Lequel ?
\[
  f_1(x)=\frac1x \qquad f_2(x)=\frac{x+5}{x-2} \qquad f_3(x)=\sqrt{x}
\]

\exo[Valeurs interdites]
Pour chaque fonction $f$, déterminer la valeur interdite.
\begin{multicols}{2}
\begin{enumerate}[label=\alph*)]
  \item ${f(x)=\frac{1}{x-6}}$
  \item ${f(x)=\frac{2}{x+5}}$
  \item ${f(x)=\frac{5}{3x}}$
  \item ${f(x)=\frac{x+2}{2x-8}}$
  \item ${f(x)=\frac{3}{2x+6}}$
  \item ${f(x)=\frac{x-1}{4x-12}}$
\end{enumerate}
\end{multicols}

\exo[Racines carrées]
Pour chaque fonction $f$, déterminer l'ensemble de définition, écrit avec un intervalle.
\begin{multicols}{2}
\begin{enumerate}[label=\alph*)]
  \item ${f(x)=\sqrt{x}}$
  \item ${f(x)=\sqrt{x-2}}$
  \item ${f(x)=\sqrt{x+4}}$
  \item ${f(x)=\sqrt{x-5}}$
  \item ${f(x)=\sqrt{3x-9}}$
  \item ${f(x)=\sqrt{2x+8}}$
\end{enumerate}
\end{multicols}

\exo[Pour s'entraîner]
Pour chaque fonction $f$, donner l'ensemble de définition, écrit avec des intervalles.
\begin{multicols}{2}
\begin{enumerate}[label=\alph*)]
  \item ${f(x)=3x^{2}-2}$
  \item ${f(x)=5-2x}$
  \item ${f(x)=\frac{3}{x-1}+x}$
  \item ${f(x)=\frac{2x+3}{5x+15}}$
  \item ${f(x)=\sqrt{4x-8}+1}$
\end{enumerate}
\end{multicols}

\partiefiche{Tableau de valeurs}

\exo[Lire un tableau]
Voici un tableau de valeurs d'une fonction $f$ :
\begin{center}
\renewcommand{\arraystretch}{1.25}
\begin{tabular}{|c|*{5}{c|}}
\hline
$x$ & $-2$ & $-1$ & $0$ & $1$ & $2$\\\hline
$f(x)$ & $2$ & $0$ & $-3$ & $-1$ & $4$\\\hline
\end{tabular}
\end{center}
\begin{enumerate}[label=\arabic*)]
  \item Quelle est l'image par $f$ de :\choix{$-1$}{$1$}{$-2$}
  \item Le nombre $2$ est-il un antécédent de $-1$ par $f$ ? Quel antécédent de $-1$ lit-on
  dans le tableau ?
\end{enumerate}

\exo[Lire un tableau (suite)]
Voici un tableau de valeurs d'une fonction $g$ :
\begin{center}
\renewcommand{\arraystretch}{1.25}
\begin{tabular}{|c|*{5}{c|}}
\hline
$x$ & $-3$ & $-1$ & $2$ & $4$ & $6$\\\hline
$g(x)$ & $4$ & $-2$ & $0$ & $5$ & $4$\\\hline
\end{tabular}
\end{center}
\begin{enumerate}[label=\arabic*)]
  \item Quelle est l'image par $g$ de :\choix{$4$}{$2$}{$-1$}
  \item Donner les antécédents lus dans le tableau de :\choix{$5$}{$-2$}{$4$}
\end{enumerate}

\exo[Construire un tableau]
On considère la fonction $f$ définie sur $\R$ par ${f(x)=0{,}5(x-1)^{2}-2}$.
\begin{enumerate}[label=\arabic*)]
  \item Dresser le tableau de valeurs de $f$ de $-2$ à $4$ avec un pas de $1$.
  \item Donner, d'après le tableau, deux antécédents de $0$ et deux antécédents de
  $2{,}5$.
  \item Tracer la courbe de $f$ dans un repère.
\end{enumerate}

\partiefiche{Courbe représentative}

\exo[Lecture graphique]
\begin{center}
\begin{tikzpicture}[x=0.5cm,y=0.5cm]
  \repere{-5}{-3}{6}{4}
  \gradx{-4,-3,-2,-1,1,2,3,4,5}
  \grady{-2,-1,1,2,3}
  \draw[coulCourbe,line width=1.1pt] (-4,2)
    .. controls (-3.667,2.333) and (-3.333,3) .. (-3,3)
    .. controls (-2.667,3) and (-2.333,2.5) .. (-2,2)
    .. controls (-1.667,1.5) and (-1.333,0.5) .. (-1,0)
    .. controls (-0.667,-0.5) and (-0.333,-0.667) .. (0,-1)
    .. controls (0.333,-1.333) and (0.667,-2) .. (1,-2)
    .. controls (1.333,-2) and (1.667,-1.333) .. (2,-1)
    .. controls (2.333,-0.667) and (2.667,-0.5) .. (3,0)
    .. controls (3.333,0.5) and (3.667,1.5) .. (4,2)
    .. controls (4.333,2.5) and (4.667,2.667) .. (5,3);
  \path[coulCourbe] (-4,2) node[croix]{} (5,3) node[croix]{};
  \node[coulCourbe,font=\footnotesize] at (-4.3,3.2) {$\Cf$};
\end{tikzpicture}
\end{center}
La courbe ci-dessus représente une fonction $f$. Lire :
\begin{enumerate}[label=\alph*)]
  \item l'ensemble de définition de $f$ ;
  \item l'image de $3$ et l'image de $-1$ ;
  \item $f(-2)$ et $f(-4)$ ;
  \item les antécédents de $2$, de $0$, de $-2$ et de $-3$ ;
  \item les nombres $x$ tels que ${f(x)=3}$.
\end{enumerate}

\exo[Points et courbes]
\begin{enumerate}[label=\arabic*)]
  \item $\Cf$ est la courbe d'une fonction $f$. Traduire chaque égalité par une phrase
  \og le point \dots{} appartient à $\Cf$\fg{} :\choix{$f(0)=-3$}{$f(-4)=0$}{$f(2)=f(6)=1$}
  \item Traduire par des égalités : a) $\Cf$ passe par le point $(3\,;-2)$ ; b) $\Cf$
  coupe l'axe des ordonnées au point d'ordonnée $4$ ; c) $\Cf$ coupe l'axe des abscisses
  aux points d'abscisses $-1$ et $5$.
  \item Pour chaque fonction $f$, dire si les points $A$, $B$ et $C$ sont sur sa courbe,
  puis donner les coordonnées de deux autres points de cette courbe.
  \begin{enumerate}[label=\alph*)]
    \item ${f(x)=x^{2}-3}$ : \ $A(-2\,;1)$, \ $B(3\,;5)$, \ $C(1{,}5\,;-0{,}75)$.
    \item ${f(x)=3x+1}$ : \ $A(-2\,;-5)$, \ $B(2\,;6)$, \ $C\!\left(\frac13\,;2\right)$.
  \end{enumerate}
\end{enumerate}

\exo[Trois antécédents]
\begin{minipage}[t]{0.52\linewidth}
\vspace{0pt}
\raggedright
À l'aide de la courbe de la fonction $f$ ci-contre, déterminer :
\begin{enumerate}[label=\alph*)]
  \item l'ensemble des nombres qui ont une image ;
  \item l'image de $2$ ;
  \item les antécédents de $-3$ ;
  \item un nombre qui a trois antécédents ;
  \item un nombre qui n'a aucun antécédent.
\end{enumerate}
\end{minipage}\hfill
\begin{minipage}[t]{0.45\linewidth}
\vspace{0pt}
\centering
\begin{tikzpicture}[x=0.4cm,y=0.4cm]
  \repere{-4}{-4}{4}{5}
  \gradx{-3,-2,-1,1,2,3}
  \grady{-3,-2,-1,1,2,3,4}
  \draw[coulCourbe,line width=1.1pt] (-3,-3)
    .. controls (-2.667,-2) and (-2.333,-0.833) .. (-2,0)
    .. controls (-1.667,0.833) and (-1.333,2) .. (-1,2)
    .. controls (-0.667,2) and (-0.333,0.5) .. (0,0)
    .. controls (0.333,-0.5) and (0.667,-1) .. (1,-1)
    .. controls (1.333,-1) and (1.667,-0.833) .. (2,0)
    .. controls (2.333,0.833) and (2.667,2.667) .. (3,4);
  \path[coulCourbe] (-3,-3) node[croix]{} (3,4) node[croix]{};
  \node[coulCourbe,font=\footnotesize] at (2.2,4.3) {$\Cf$};
\end{tikzpicture}
\end{minipage}

\exo[Quelle courbe passe par ce point ?]
On considère les fonctions définies sur $\intoo{0}{+\infty}$ par :
\[
  f(x)=10-x \quad ; \qquad g(x)=\frac{24}{x} \quad ;
\]
\[
  h(x)=-2x+14 \quad ; \qquad k(x)=-\frac12x^{2}+5x-4.
\]
\begin{enumerate}[label=\arabic*)]
  \item Parmi ces fonctions, lesquelles ont une courbe qui passe par le point $A(4\,;6)$ ?
  \item Même question avec $B(6\,;4)$, puis avec $C(8\,;4)$.
\end{enumerate}

\partiefiche{Calculatrice}

\exo[Antécédents à la calculatrice]
\begin{enumerate}[label=\arabic*)]
  \item Afficher le tableau de valeurs de la fonction $f$ définie sur $\intff{-2}{6}$ par
  ${f(x)=0{,}5x^{2}-2x+1}$, de $-2$ à $6$ avec un pas de $1$, puis tracer sa courbe à
  l'écran avec des axes adaptés.
  \item Avec \texttt{Calculer}, \texttt{Antécédent}, donner les antécédents de $2$ par
  $f$ arrondis à $0{,}1$.
  \item Le nombre $0$ a-t-il un antécédent ? plusieurs ?
\end{enumerate}

\exo[Tableau avec un pas de 0,5]
Soit $g$ la fonction définie sur $\intff{-1}{2}$ par ${g(x)=-2x^{2}+x+3}$.
\begin{enumerate}[label=\arabic*)]
  \item Afficher le tableau de valeurs de $g$ de $-1$ à $2$ avec un pas de $0{,}5$, et le
  recopier.
  \item Tracer la courbe de $g$ à l'écran, puis sur la copie.
  \item Donner $g(0{,}25)$ et les antécédents de $2$ et de $0$.
\end{enumerate}

\exo[Une fonction quotient]
On considère la fonction $f$ définie sur $\intff{0}{5}$ par
\[ f(x)=\frac{2x-1}{(x+1)^{2}}. \]
\begin{enumerate}[label=\arabic*)]
  \item Dresser le tableau de valeurs de $f$ de $0$ à $5$ avec un pas de $0{,}5$ (valeurs
  arrondies à $0{,}01$).
  \item Tracer la courbe de $f$ à la calculatrice.
  \item Quels sont les antécédents de $0{,}25$ ?
  \item Résoudre graphiquement ${f(x)=0}$, puis $f(x)\ge 0{,}25$.
  \item Tracer aussi la droite d'équation $y=x-1$, puis résoudre graphiquement
  $\dfrac{2x-1}{(x+1)^{2}}=x-1$ (valeurs arrondies à $0{,}01$).
\end{enumerate}

\exo[Intersections de courbes]
\begin{enumerate}[label=\arabic*)]
  \item Tracer à l'écran les courbes des fonctions $f$ et $g$ définies sur
  $\intff{-2}{3}$ par ${f(x)=x^{3}-2x^{2}}$ et ${g(x)=x-2}$, puis résoudre
  graphiquement $x^{3}-2x^{2}\le x-2$.
  \item Trouver les coordonnées du ou des points communs aux courbes d'équations
  $y=3x^{2}+x+5$ et $y=3x^{2}-2x+8$. Confirmer par le calcul.
  \item Même question avec $y=\dfrac1x$ et $y=\dfrac{1+2x}{x}$.
\end{enumerate}

\partiefiche{Résolution graphique}

\exo[Équations]
Ci-dessous, la courbe d'une fonction $f$ définie sur $\intff{0}{10}$.
\begin{center}
\begin{tikzpicture}[x=0.5cm,y=0.5cm]
  \repere{-1}{-3}{11}{3}
  \gradx{1,2,3,4,5,6,7,8,9,10}
  \grady{-2,-1,1,2}
  \draw[coulCourbe,line width=1.1pt] (0,0)
    .. controls (0.333,-0.667) and (0.667,-2) .. (1,-2)
    .. controls (1.333,-2) and (1.667,-1.333) .. (2,-1)
    .. controls (2.333,-0.667) and (2.667,-0.333) .. (3,0)
    .. controls (3.333,0.333) and (3.667,0.667) .. (4,1)
    .. controls (4.333,1.333) and (4.667,2) .. (5,2)
    .. controls (5.333,2) and (5.667,1.222) .. (6,1)
    .. controls (6.667,0.556) and (7.333,0.333) .. (8,0)
    .. controls (8.667,-0.333) and (9.333,-0.667) .. (10,-1);
  \path[coulCourbe] (0,0) node[croix]{} (10,-1) node[croix]{};
  \node[coulCourbe,font=\footnotesize] at (5,2.5) {$\Cf$};
\end{tikzpicture}
\end{center}
\begin{enumerate}[label=\arabic*)]
  \item Résoudre graphiquement :
  \choix{${f(x)=2}$}{${f(x)=1}$}{${f(x)=0}$}
  \par\hspace*{1em}\makebox[0.3\linewidth][l]{d) ${f(x)=-2}$}e) ${f(x)=3}$
  \item Combien de solutions l'équation ${f(x)=-1}$ a-t-elle ?
\end{enumerate}

\exo[Inéquations]
La courbe de $f$ est celle de l'exercice précédent. Résoudre graphiquement :
\begin{multicols}{2}
\begin{enumerate}[label=\alph*)]
  \item $f(x)>-2$
  \item $f(x)<1$
  \item $f(x)\le 2$
  \item $f(x)\ge 1$
\end{enumerate}
\end{multicols}

\exo[Équation et inéquations]
\begin{minipage}[t]{0.50\linewidth}
\vspace{0pt}
\raggedright
$\mathcal{C}$ est la courbe d'une fonction $f$ définie sur $\intff{-4}{5}$.
\begin{enumerate}[label=\alph*)]
  \item Résoudre graphiquement ${f(x)=1}$.
  \item Résoudre graphiquement $f(x)\le 1$.
  \item Résoudre graphiquement $f(x)<-1$.
\end{enumerate}
\end{minipage}\hfill
\begin{minipage}[t]{0.48\linewidth}
\vspace{0pt}
\centering
\begin{tikzpicture}[x=0.38cm,y=0.38cm]
  \repere{-5}{-3}{6}{4}
  \gradx{-4,-3,-2,-1,1,2,3,4,5}
  \grady{-2,-1,1,2,3}
  \draw[coulCourbe,line width=1.1pt] (-4,1)
    .. controls (-3.667,0.333) and (-3.333,-0.5) .. (-3,-1)
    .. controls (-2.667,-1.5) and (-2.333,-2) .. (-2,-2)
    .. controls (-1.667,-2) and (-1.333,-1.5) .. (-1,-1)
    .. controls (-0.667,-0.5) and (-0.333,0.5) .. (0,1)
    .. controls (0.333,1.5) and (0.667,1.667) .. (1,2)
    .. controls (1.333,2.333) and (1.667,3) .. (2,3)
    .. controls (2.333,3) and (2.667,2.333) .. (3,2)
    .. controls (3.333,1.667) and (3.667,1.333) .. (4,1)
    .. controls (4.333,0.667) and (4.667,0.333) .. (5,0);
  \path[coulCourbe] (-4,1) node[croix]{} (5,0) node[croix]{};
  \node[coulCourbe,font=\footnotesize] at (1,3.4) {$\mathcal{C}$};
\end{tikzpicture}
\end{minipage}

\exo[Avec la formule]
Soit $f$ la fonction définie sur $\intff{-3}{3}$ par
${f(x)=\frac{x^{3}-7x+6}{2}}$, représentée ci-dessous.
\begin{center}
\begin{tikzpicture}[x=0.5cm,y=0.5cm]
  \repere{-4}{-2}{4}{8}
  \gradx{-3,-2,-1,1,2,3}
  \grady{-1,1,2,3,4,5,6,7}
  \draw[coulCourbe,line width=1.1pt,domain=-3:3,samples=80] plot (\x,{((\x)*(\x)*(\x)-7*(\x)+6)/2});
  \node[coulCourbe,font=\footnotesize] at (2.3,6) {$\Cf$};
\end{tikzpicture}
\end{center}
\begin{enumerate}[label=\arabic*)]
  \item Dresser à la calculatrice le tableau de valeurs de $f$ de $-3$ à $3$ avec un pas
  de $1$.
  \item Lire l'image de $0$, puis résoudre graphiquement ${f(x)=6}$.
  \item Vérifier que $x^{3}-7x-6=(x+1)(x+2)(x-3)$ et en déduire les solutions exactes
  de ${f(x)=6}$.
  \item Résoudre graphiquement $f(x)\ge 6$, puis $f(x)<0$.
  \item Combien de solutions ont les équations ${f(x)=0}$ ; ${f(x)=4}$ ; ${f(x)=-1}$ ?
\end{enumerate}

\partiefiche{Comparer deux fonctions}

\exo[Deux courbes]
Les courbes ci-dessous représentent deux fonctions $f$ et $g$ définies sur
$\intff{-3}{7}$.
\begin{center}
\begin{tikzpicture}[x=0.5cm,y=0.5cm]
  \repere{-4}{-2}{8}{4}
  \gradx{-3,-2,-1,1,2,3,4,5,6,7}
  \grady{-1,1,2,3}
  \draw[coulCourbe,line width=1.1pt] (-3,1)
    .. controls (-2.667,1.333) and (-2.333,1.667) .. (-2,2)
    .. controls (-1.667,2.333) and (-1.333,3) .. (-1,3)
    .. controls (-0.667,3) and (-0.333,2.333) .. (0,2)
    .. controls (0.333,1.667) and (0.667,1.333) .. (1,1)
    .. controls (1.333,0.667) and (1.667,0.333) .. (2,0)
    .. controls (2.333,-0.333) and (2.667,-1) .. (3,-1)
    .. controls (3.333,-1) and (3.667,-0.333) .. (4,0)
    .. controls (4.333,0.333) and (4.667,0.667) .. (5,1)
    .. controls (5.333,1.333) and (5.667,1.667) .. (6,2)
    .. controls (6.333,2.333) and (6.667,2.667) .. (7,3);
  \draw[coulCourbeC,line width=1.1pt] (-3,3)
    .. controls (-2.667,2.667) and (-2.333,2.333) .. (-2,2)
    .. controls (-1.667,1.667) and (-1.333,1.333) .. (-1,1)
    .. controls (-0.667,0.667) and (-0.333,0.333) .. (0,0)
    .. controls (0.333,-0.333) and (0.667,-1) .. (1,-1)
    .. controls (1.333,-1) and (1.667,-0.333) .. (2,0)
    .. controls (2.333,0.333) and (2.667,0.667) .. (3,1)
    .. controls (3.333,1.333) and (3.667,1.667) .. (4,2)
    .. controls (4.333,2.333) and (4.667,3) .. (5,3)
    .. controls (5.333,3) and (5.667,2.333) .. (6,2)
    .. controls (6.333,1.667) and (6.667,1.333) .. (7,1);
  \node[coulCourbe,font=\footnotesize,right] at (7,3) {$\Cf$};
  \node[coulCourbeC,font=\footnotesize,right] at (7,1) {$\Cg$};
\end{tikzpicture}
\end{center}
Résoudre graphiquement :
\begin{enumerate}[label=\alph*)]
  \item l'équation ${f(x)=g(x)}$ ;
  \item l'inéquation $f(x)>g(x)$ ;
  \item l'inéquation $f(x)\le g(x)$.
\end{enumerate}

\exo[Une droite et une courbe]
$\Cf$ (droite) et $\Cg$ représentent deux fonctions $f$ et $g$ définies sur
$\intff{-6}{6}$.
\begin{center}
\begin{tikzpicture}[x=0.45cm,y=0.45cm]
  \repere{-7}{-3}{7}{6}
  \gradx{-6,-5,-4,-3,-2,-1,1,2,3,4,5,6}
  \grady{-2,-1,1,2,3,4,5}
  \draw[coulCourbeC,line width=1.1pt] (-6,5)--(6,-1);
  \draw[coulCourbe,line width=1.1pt] (-6,1)
    .. controls (-5.667,1.667) and (-5.333,2.5) .. (-5,3)
    .. controls (-4.667,3.5) and (-4.333,3.667) .. (-4,4)
    .. controls (-3.667,4.333) and (-3.333,5) .. (-3,5)
    .. controls (-2.667,5) and (-2.333,4.222) .. (-2,4)
    .. controls (-1.333,3.556) and (-0.667,3.444) .. (0,3)
    .. controls (0.333,2.778) and (0.667,2.333) .. (1,2)
    .. controls (1.333,1.667) and (1.667,1.333) .. (2,1)
    .. controls (2.333,0.667) and (2.667,0.333) .. (3,0)
    .. controls (3.333,-0.333) and (3.667,-0.667) .. (4,-1)
    .. controls (4.333,-1.333) and (4.667,-2) .. (5,-2)
    .. controls (5.333,-2) and (5.667,-1.333) .. (6,-1);
  \node[coulCourbeC,font=\footnotesize,above right] at (-6,5) {$\Cf$};
  \node[coulCourbe,font=\footnotesize,left] at (-6,1) {$\Cg$};
\end{tikzpicture}
\end{center}
Résoudre graphiquement :
\begin{multicols}{2}
\begin{enumerate}[label=\alph*)]
  \item $f(x)\le 2$
  \item $g(x)<0$
  \item $g(x)\ge 4$
  \item ${f(x)=g(x)}$
  \item $f(x)<g(x)$
\end{enumerate}
\end{multicols}

\exo[Graphique et calcul]
On considère les fonctions $f$ et $g$ définies sur $\intff{-2}{3}$ par
${f(x)=x^{2}-2}$ et ${g(x)=x}$.
\begin{enumerate}[label=\arabic*)]
  \item Dresser les tableaux de valeurs de $f$ et de $g$ de $-2$ à $3$ avec un pas de $1$.
  \item Tracer les deux courbes dans un même repère.
  \item Résoudre graphiquement ${f(x)=g(x)}$, puis $f(x)<g(x)$.
  \item Vérifier que $x^{2}-2-x=(x+1)(x-2)$, et retrouver ainsi les solutions de
  ${f(x)=g(x)}$.
  \item Le point $P(1{,}5\,;0{,}25)$ est-il sur la courbe de $f$ ? sur celle de $g$ ?
\end{enumerate}

\end{multicols}

\end{document}
